% --- scholar LaTeX fragment --------------------------------------------
% Generated by: scholar sessions export -f table
% Session: classes-fenomenografi
% \input this file from the body of a document whose
% preamble loads:
%   \usepackage{longtable}
%   \usepackage{booktabs}
%   \usepackage{array}
% ----------------------------------------------------------------------

\begingroup\footnotesize
\begin{longtable}{@{}%
  >{\raggedright\arraybackslash}p{0.44\textwidth}c%
  >{\raggedright\arraybackslash}p{0.13\textwidth}%
  >{\raggedright\arraybackslash}p{0.26\textwidth}@{}}
\caption{Träffar som rör påståendet, uppfattningar av klass och objekt (24 citerade, 22 som stöder påståendet och 6 som kvalificerar eller motsäger det, av 649 unika träffar; de 161 angränsande och 436 felträffarna listas inte). Skälet till varje inkludering står i sista kolumnen; för maskinklassade rader anges modellens konfidens. Databaser: OpenAlex (OA; inte open access), IEEE Xplore (IEEE), Web of Science (WoS), Scopus.}\label{tab:sok-fenomenografi}\\
\toprule
Titel & År & Leverantör & Skäl \\
\midrule
\endfirsthead
\multicolumn{4}{@{}l}{\emph{\tablename~\thetable{} (forts.)}}\\
\toprule
Titel & År & Leverantör & Skäl \\
\midrule
\endhead
\midrule \multicolumn{4}{r@{}}{\emph{forts.\ på nästa sida}}\\
\endfoot
\bottomrule
\endlastfoot
A Systematic Review of Threshold Concept Identification in Programming Education: Foundations for MOOC-Based Learning Analytics & 2027 & Scopus & \emph{citerad}: tröskelbegrepp \\
Difficulties in Learning Inheritance and Polymorphism & 2011 & OA & \emph{citerad}: verklighetsmodellen \\
From Phenomenography Study to Planning Teaching & 2010 & WoS, Scopus & \emph{citerad}: undervisning \\
Fuzzy OOP: Expanded and reduced term interpretations & 2012 & Scopus & \emph{citerad}: nivåer av klass och objekt \\
Novice Java programmers' conceptions of \textquotedbl{}object\textquotedbl{} and \textquotedbl{}class\textquotedbl{}, and variation theory & 2005 & OA, Scopus & \emph{citerad}: nivåer av klass och objekt \\
Novice students' learning of object-oriented programming & 2006 & OA & \emph{citerad}: nivåer av klass och objekt \\
Object-oriented design and programming: An investigation of novices' conceptions on objects and classes & 2015 & Scopus & \emph{citerad}: nivåer av klass och objekt \\
Partonomy and taxonomy in object-oriented thinking & 2006 & OA & \emph{citerad}: verklighetsmodellen \\
Phenomenography: From Critical Aspects to Knowledge Claim & 2013 & OA & \emph{citerad}: ansatsen och fyndens hållbarhet \\
Programming Course Design: Phenomenographic Approach to Learning and Teaching & 2014 & IEEE & \emph{citerad}: undervisning \\
Reflections on threshold concepts in computer programming and beyond & 2010 & OA & \emph{citerad}: tröskelbegrepp \\
Relationship between text and action conceptions of programming: A phenomenographic and quantitative perspective & 2011 & Scopus & \emph{citerad}: text och handling \\
Statistical frequency-analysis of misconceptions in object-oriented-programming - Regularized PCR models for frequency analysis across OOP concepts and related factors & 2018 & Scopus & \emph{citerad}: ansatsen och fyndens hållbarhet \\
Student understanding of object-oriented programming as expressed in concept maps & 2008 & OA & \emph{citerad}: nivåer av klass och objekt \\
Students' different understandings of class diagrams & 2012 & Scopus & \emph{citerad}: relationer mellan klasser \\
Students' understandings of storing objects & 2007 & OA & \emph{citerad}: lagring och identitet \\
Teaching Object-Oriented Threshold Concepts Using the Learning-from-Errors Approach & 2025 & Scopus & \emph{citerad}: tröskelbegrepp \\
THE EFFECTS OF INTEGRATING VARIATION THEORY ON A WEB-BASED OBJECT-ORIENTED PROGRAMMING LEARNING & 2012 & WoS & \emph{citerad}: undervisning \\
The nature of an object-oriented program: How do practitioners understand the nature of what they are creating? & 2011 & Scopus & \emph{citerad}: relationer mellan klasser \\
The same but different students' understandings of primitive and object variables & 2008 & OA & \emph{citerad}: lagring och identitet \\
Threshold concepts in computer science & 2007 & OA & \emph{citerad}: tröskelbegrepp \\
Using variation theory as a guiding principle in an OOP assisted syntax correction learning system & 2020 & Scopus & \emph{citerad}: undervisning \\
What makes the difference? An empirical comparison of critical aspects identified in phenomenographic and variation theory analyses & 2019 & OA & \emph{citerad}: ansatsen och fyndens hållbarhet \\
When intuition and logic clash: The case of the object-oriented paradigm & 2012 & OA & \emph{citerad}: verklighetsmodellen \\
A cognitive model of how interactive multimedia,authoring facilitates conceptual understanding of object-Oriented programming in novices & 2011 & Scopus & stöder påståendet (0.87) \\
A Long-Term Investigation of the Comprehension of OOP Concepts by Novices & 2005 & WoS & stöder påståendet (1.00) \\
Conceptions of the students around object-oriented design: A case study & 2017 & Scopus & stöder påståendet (0.89) \\
Development of object-understanding among students in the humanities & 2005 & Scopus & stöder påståendet (0.98) \\
Encapsulation and reuse as viewed by Java students & 2001 & WoS & stöder påståendet (0.99) \\
Identifying the difficulties of object-oriented development & 2002 & WoS & stöder påståendet (0.95) \\
Internal representation and rule development in object-oriented design & 1995 & OA & stöder påståendet (0.93) \\
Novice Comprehension of Object-Oriented OO Programs: An Empirical Study & 2015 & WoS, IEEE & stöder påståendet (0.89) \\
Object oriented programming and program correctness: The students' perspective & 2006 & Scopus & stöder påståendet (0.96) \\
Object-Oriented Program Comprehension: Effect of Expertise, Task and Phase & 2002 & OA & stöder påståendet (0.95) \\
Object-oriented programming - What do students think of objects and classes? & 2011 & Scopus & stöder påståendet (0.97) \\
Object-Oriented Programming: Diagnosis Understanding by Identifying and Describing Novice Perceptions & 2020 & IEEE & stöder påståendet (0.90) \\
Pre-service and in-service teachers' experiences of learning to program in an object-oriented language & 2008 & Scopus & stöder påståendet (0.98) \\
Representation of abstract concepts: Differences across computing disciplines & 2015 & IEEE & stöder påståendet (0.74) \\
Solving introductory computer programming problems using an object-oriented language: Insights from an empirical study & 2008 & Scopus & stöder påståendet (0.94) \\
Students' Conceptions of Programming in the Context of Game Design & 2022 & Scopus & stöder påståendet (0.98) \\
Teaching Sequence Diagrams To Programming Beginners And The Change Of Algorithmic Conceptions & 2017 & WoS, IEEE & stöder påståendet (0.92) \\
Thinking in imperative or objects? A study on how novice programmer thinks when it comes to designing an application & 2019 & IEEE & stöder påståendet (0.91) \\
Understanding the ``this'' reference in object oriented programming: Misconceptions, conceptions, and teaching recommendations & 2021 & Scopus & stöder påståendet (1.00) \\
Variation in students' conceptions of object-oriented information system development & 2005 & WoS & stöder påståendet (0.97) \\
What does it take to learn 'programming thinking'? & 2005 & Scopus & stöder påståendet (0.94) \\
What students (should) know about object oriented programming & 2011 & OA & stöder påståendet (0.97) \\
A comparison of the comprehension of object-oriented and procedural programs by novice programmers & 1999 & IEEE & kvalificerar eller motsäger påståendet (0.98) \\
An empirical study of novice program comprehension in the imperative and object-oriented styles & 1997 & OA & kvalificerar eller motsäger påståendet (0.99) \\
Assessing the cognitive consequences of the object-oriented approach: A survey of empirical research on object-oriented design by individuals and teams & 1997 & OA & kvalificerar eller motsäger påståendet (0.95) \\
Is there any difference in novice comprehension of a small program written in the event-driven and object-oriented styles? & 2002 & WoS & kvalificerar eller motsäger påståendet (0.95) \\
Research perspectives on the objects-early debate & 2006 & OA & kvalificerar eller motsäger påståendet (0.94) \\
The impact of experience on individual performance and workload differences using object-oriented and process-oriented systems analysis techniques & 1996 & IEEE & kvalificerar eller motsäger påståendet (0.98) \\
\end{longtable}
\endgroup
